\chapter{Conclusion}

It has been clearly shown that job seekers face a very diverse range of problems which inhibit them from finding jobs or starting businesses.

Job seekers sometimes fail to find jobs because employers seek years of experience that they do not have (and sometimes are not willing to train unskilled workers).

From the industries’ perspectives: the industries’ skills requirements do not match the content taught at tertiary institutions, there is a shortage of jobs, some jobs require high proficiency in English and Afrikaans which some struggle to speak because their primary education was offered in indigenous languages like IsiXhosa and isiZulu.

Job seekers are unaware of where to find job vacancies, lack funds to study to improve their skills and gain further qualifications, lack nationally recognised qualifications which are required for several jobs, job seekers do not have IDs whereas several South African systems are built to work only with ID numbers and do not work with other forms of identification like asylum passports and other foreign passports. Job seekers lack nationally recognised certificates for skills although they sometimes have the skills thus we recommend the usage of cover letters, and social media referrals for the low skilled to assist them in securing short-term jobs and gaining trust which in turn gives them an opportunity to gain experience.

During the research period, a number of job seekers indicated that if they were provided with capital they would start a business and if they received funding for schooling they would attend school. Lack of capital to start businesses and lack of study loans for job seekers can be provided by governmental funding associations, NGOs and private funding associations.

However, tasks that previously required human labour such as broadcasting vacancies, creating CVs, seeking for jobs door-to-door can be mitigated by web scrappers, job finding portals, online training apps and other specialised applications. Moreover, developmental ICTs can even be used to provide jobseekers with information about where to find the resources they need. Yet certain tasks still require human assistance, thus, the need to work with intermediaries with better ICT skills than the low-skilled job seekers. Alternatively, job seekers can be paired up when using complex ICT based platforms. On the contrary, designers can avoid the need for assistance by developing applications to be easy to learn and easy to use.

We found out that education and training for ICT usage are very vital for many ICT interventions to achieve their intended functions as unemployment eradication tools. The lack of technological know-how prevents low-skilled individuals from efficiently utilising generic resources which are used by the public, for example existing job listing platforms such as Career24 and Indeed.

In contexts where participants are not available throughout the research period, the usage of Co-design is challenged and the specific methods for engaging the co-designers should be carefully thought out. Larger setups like workshops and focus groups work to gain multiple inputs within a very limited period whereas one-on-one methods better serve to get users’ uninfluenced preferences and values, which make the findings of research more trustworthy. Also, intermediary input can be used when researchers have limited time with co-designers. However, the information obtained from intermediaries as experts should not reduce the attention given to first-hand information as well as observations from the individuals concerned (in this case job seekers) as researchers might miss major findings while focusing on knowledge based on experience or expertise.

We have shown that systems that provide information that is already filtered for their users like Khanyo Job Network work well in dealing with low-skilled individuals, to bring their ICT benefits to par with the benefits available to individuals with high technological know-how. Overall, we found that Khanyo Job Network is promising as all the participants were willing to use it, and in its two weeks of deployment, one of the participants had gotten a job through the application.